\documentclass[12pt]{article}

\usepackage{jheppub}
\usepackage{iip-arxiv}
\toccontinuoustrue % Não quebra a página título
\compress % Comprime o resumo e sumário para caber em uma página

\usepackage{amsmath}
\usepackage{amssymb}
\usepackage{braket}
\usepackage{xcolor}

\title{Notes on VQE for Quantum Chemistry}
\author[a]{Fábio Novaes}
\affiliation[a]{Federal Rural University of Pernambuco}
\emailAdd{fabio.nsantos@gmail.com}
\date{\today}
\abstract{}

\begin{document}
\maketitle

\section{Introduction}

Classical computational chemistry methods, such as Hartree-Fock and Density Functional Theory, often struggle to accurately describe strongly correlated systems. The Variational Quantum Eigensolver (VQE) is a promising quantum algorithm that can potentially overcome these limitations by leveraging the power of quantum computers to find the ground state energy of molecular systems. In this set of notes, we will explore the VQE algorithm in the context of quantum chemistry, discussing its implementation, challenges, and potential applications.

One of the first papers to suggest the use of quantum computers for quantum chemistry was by Aspuru-Guzik et al.~in 2005~\cite{aspuru-guzikSimulatedQuantumComputation2005}. They proposed using a quantum computer to simulate the electronic structure of molecules, which is a fundamental problem in chemistry. The VQE algorithm was later developed as a hybrid quantum-classical approach to find the ground state energy of molecular systems efficiently \cite{peruzzo2014variational,mcclean2016theory}. A review of the VQE algorithm and its applications in quantum chemistry can be found in~\cite{cerezoVariationalQuantumAlgorithms2021}. 

\section{How to Obtain Molecular Hamiltonians?}

To apply the VQE algorithm to quantum chemistry problems, we first need to obtain the molecular Hamiltonian in a form that can be implemented on a quantum computer. This typically involves several steps, including choosing a basis set, performing a Hartree-Fock calculation, and mapping the resulting fermionic Hamiltonian to a qubit Hamiltonian. The choice of basis set can significantly affect the accuracy of the results, with common choices including STO-3G, 6-31G, and cc-pVDZ. The Hartree-Fock calculation provides the molecular orbitals and their corresponding energies, which are used to construct the fermionic Hamiltonian. Finally, the fermionic Hamiltonian is mapped to a qubit Hamiltonian using techniques such as the Jordan-Wigner or Bravyi-Kitaev transformations. The resulting qubit Hamiltonian can then be used as the input for the VQE algorithm, where we will optimize a parameterized quantum circuit to find the ground state energy of the molecule. 

Let us consider as examples the following molecules: H$_2$, LiH, and BeH$_2$. For each molecule, we will discuss the choice of basis set, the resulting Hamiltonian, and the implementation of the VQE algorithm to find the ground state energy.

Following \cite{kandalaHardwareefficientVariationalQuantum2017}, we choose the STO-3G basis set, which corresponds to the functional form
\begin{equation}
    \chi_i(\mathbf{r}) = \sum_{j=1}^3 c_{ij} \phi_{ij}(\mathbf{r}),
\end{equation}
where $\phi_{ij}(\mathbf{r})$ are Gaussian functions and $c_{ij}$ are coefficients that are optimized to best fit the true atomic orbitals labeled by $i$. For the H$_2$ molecule, each atom contributes a 1s orbital, for a total of 4 spin-orbitals. We set the X axis as the interatomic axis for the LiH and BeH$_2$ molecules, and consider the orbitals 1s for each H atom and 1s, 2s, 2p$_x$ for the Li and Be atoms, assuming zero filling for the 2p$_y$ and 2p$_z$ orbitals, which do not interact strongly with the subset of orbitals considered \cite{kandalaHardwareefficientVariationalQuantum2017}. 

The second quantized Hamiltonian can be expressed in the form
\begin{equation}
    H = H_1 + H_2 = \sum_{p,q=1}^M h_{pq} a_p^\dagger a_q + \frac{1}{2} \sum_{p,q,r,s=1}^M h_{pqrs} a_p^\dagger a_q^\dagger a_r a_s,
\end{equation}
where $a_p^\dagger$ and $a_q$ are fermionic creation and annihilation operators, and $h_{pq}$ and $h_{pqrs}$ are the one-electron and two-electron integrals, respectively. Here $M=4, 8, 10$
is the number of spin-orbitals for H$_2$, LiH, and BeH$_2$, respectively. The coefficients $h_{pq}$ and $h_{pqrs}$ are computed using the chosen basis set and the molecular geometry, and they encode the interactions between the electrons in the molecule as
\begin{equation}
    h_{pq} = \int d\mathbf{r} \, \chi_p^*(\mathbf{r}) \left( -\frac{1}{2} \nabla^2 - \sum_A \frac{Z_A}{|\mathbf{r} - \mathbf{R}_A|} \right) \chi_q(\mathbf{r}),
\end{equation}
where $Z_A$ is the atomic number of nucleus $A$ located at $\mathbf{R}_A$. The two-electron integrals $h_{pqrs}$ represent the Coulomb repulsion between pairs of electrons and are given by
\begin{equation}
    h_{pqrs} = \int d\mathbf{r}_1 d\mathbf{r}_2 \, \frac{\chi_p^*(\mathbf{r}_1) \chi_q^*(\mathbf{r}_2) \chi_r(\mathbf{r}_1) \chi_s(\mathbf{r}_2)}{|\mathbf{r}_1 - \mathbf{r}_2|},
\end{equation} 
where $\mathbf{r}_1$ and $\mathbf{r}_2$ are the coordinates of the two electrons. It is assumed that the spin is conserved, so the integrals are zero unless the spin of the orbitals $p$ and $r$ are the same, and the spin of the orbitals $q$ and $s$ are the same.


\bibliographystyle{JHEP}
\bibliography{qth-chemistry-vqe}

\end{document}
